% ============================================================================================
\section{I3: Horizon and Ridgeline Plots (Owner: Shrey Modi, BT2024125)}
\label{sec:i3}
Both techniques are implemented by hand in Matplotlib (\fpath{infovis/i3_horizon_ridgeline/run.py}); no horizon or ridgeline
library is used.

\subsection{Horizon plot construction}
Horizon graphs~\cite{saito2005,heer2009} fold a time series into a compact band:
\begin{enumerate}
\item centre the series on a baseline $b$ (the row's mean or median): $d_t=x_t-b$;
\item choose $N$ bands of height $h=\max|d|/N$, shared by all rows of a figure so that rows are comparable;
\item slice the positive part into bands $p_t^{(k)}=\mathrm{clip}(d_t-kh,0,h)$ and the negative part into
  $n_t^{(k)}=\mathrm{clip}(-d_t-kh,0,h)$, for $k=0..N{-}1$;
\item \emph{fold} the negative bands upward onto the same $[0,h]$ strip and overlay all bands, drawing band $k$ in the
  $k$-th saturation of a diverging palette (reds above the baseline, blues below; RdBu).
\end{enumerate}
Fig.~\ref{fig:35} shows each step on the monthly ``Q\&A and advice'' series and compares $N=1,2,3,4,6$.

\begin{algorithm}[!t]
\caption{Horizon band folding for one row}
\label{alg:horizon}
\begin{algorithmic}[1]
\REQUIRE series $x_{1..T}$, baseline $b$, band height $h$, bands $N$, palettes $C^+_{1..N}$, $C^-_{1..N}$
\STATE $d_t\leftarrow x_t-b$
\FOR{$k=0$ \TO $N-1$}
  \STATE $p_t\leftarrow\min(\max(d_t-kh,0),h)$; \ fill area $[0,p_t]$ with $C^+_{k+1}$
  \STATE $n_t\leftarrow\min(\max(-d_t-kh,0),h)$; \ fill area $[0,n_t]$ with $C^-_{k+1}$
\ENDFOR
\STATE set $y$-limits to $[0,h]$ \COMMENT{values beyond $Nh$ saturate in the darkest band}
\end{algorithmic}
\end{algorithm}

\figfull{Fig_35_horizon_construction.png}{Horizon construction for ``Q\&A and advice'' monthly top-post counts. (a) Raw
series and mean baseline; (b) deviation split into 3 bands; (c) negative bands mirrored up and all bands collapsed onto
$[0,h]$; (d--h) the same series with $N=1,2,3,4,6$ bands at the same row height; (i) rationale for the band count.}{fig:35}

\textbf{Band count (Fig.~\ref{fig:35}d--h).} With $N=1$ the row is a two-colour area chart squeezed to one band: the
2019--20 surge flattens and small dips vanish. With $N=2$--3, each band covers a third or half of the range, so the
peak of September~2020 (238 posts against a mean of 72.7) reads as dark red and the 2023--24 deficit as mid blue. With
$N\ge4$ the vertical resolution grows, but four or more shades of one hue are hard to tell apart, and Heer et al.\ found that
estimation accuracy drops as the number of layers increases~\cite{heer2009}. We use $N=3$ throughout.

\subsection{Variation 1: monthly top posts per topic}
\figfull{Fig_36_horizon_topics_monthly.png}{Horizon plot of monthly all-time-top posts for the eight topic categories, Jan
2016--Aug 2024 (104 months), baseline = row mean, $h=55.1$ posts/month, $N=3$. Dotted lines: Mar~2020 (COVID-19 lockdowns)
and Jun~2023 (Reddit API protest).}{fig:36}

\textbf{Inferences.}
(1) Every topic follows the same skeleton: blue (below mean) in 2016, red (above mean) between about 2017 and early 2021, and
blue again from mid-2021. This is the \emph{survivorship} structure of an all-time-top corpus, not a change in Reddit
activity. Older posts were made when communities were smaller; the newest have had less time to collect votes.
(2) The \emph{timing} of the red phase differs by topic. Entertainment and media peaks earliest (dark red in
2017--2018; its top month is March~2018 with 185 posts). Creative and lifestyle peaks in early 2019 (113 posts in
January~2019). The discussion topics peak during the pandemic: Q\&A and advice reaches the darkest band in mid-2020 (238
posts in September~2020), News, science and tech has repeated dark peaks through 2020 to January~2021 (224 posts; the timing coincides
with the US election period, though the data alone cannot show the cause), and Memes, Photos and Text humour all peak in
July--August~2020.
(3) The March~2020 line marks the start of the strongest red period for Q\&A, Memes and Photos. That is consistent with
higher time spent on Reddit during lockdowns producing more posts that later counted as all-time tops.
(4) After the June~2023 API-protest line, Entertainment, News and Photos drop into the mid-blue band. Their recent output is
well below their own means, while Creative and lifestyle stays near its mean. The protest and blackout coincide with, but
cannot alone explain, the decline, because survivorship also penalises recent months.
(5) News has a short red spike in early 2022 that no other topic has (February--March~2022). Its timing matches the start of
the war in Ukraine, a plausible but untested explanation.
The horizon makes this visible because each row is judged against its own baseline.

\subsection{Line chart vs.\ horizon: vertical space efficiency}
\figfull{Fig_37_line_vs_horizon.png}{The same eight monthly series in equal vertical space: (a) overlapping line chart;
(b) small-multiple line charts on a shared 0--238 scale; (c) horizon rows with $h=55.1$ posts and 3 bands.}{fig:37}

\textbf{Inferences.}
(1) The overlapping line chart (Fig.~\ref{fig:37}a) shows the shared 2020 peak but hides individual topics: eight lines
cross repeatedly between 40 and 150 posts per month.
(2) Small multiples on a shared scale (Fig.~\ref{fig:37}b) separate the series, but each row has only one eighth of the
height for a 0--238 range. Low-volume topics such as Knowledge and curiosity become thin slivers.
(3) The horizon row uses the same height but only needs to resolve one band of 55.1 posts. Its effective vertical
resolution is $238/55.1\approx4.3$ times higher than the small multiple's. Above- versus below-mean phases are also coded
by hue, so they are visible pre-attentively. The trade-off is learnability: folded negatives and band shading must be
explained (the legend gives the band limits).

\subsection{Variation 2: temporal granularity (quarterly vs.\ yearly)}
\figfull{Fig_38_horizon_topics_quarter_year.png}{Granularity comparison: (a) quarterly counts, $h=150.3$ posts/quarter; (b)
yearly counts, $h=469$ posts/year (step rendering; the last bin is partial, data end 11~Sep~2024).}{fig:38}

\textbf{Inferences.}
(1) Quarterly aggregation (Fig.~\ref{fig:38}a) removes month-to-month noise while keeping turning points. Q\&A's darkest
band is now a single quarter (Q3~2020), and Entertainment's two separate peaks (2018 and 2020) remain distinct.
(2) Yearly aggregation (Fig.~\ref{fig:38}b) shows the long-run regime most clearly. 2020 is the darkest year for Q\&A,
News, Photos and Memes. Entertainment is dark red from 2017 to 2020 and Text humour in 2019--2020. All topics are blue in
2014--2016 and 2023--2024. The yearly view also loses information: it cannot show that News peaked in January~2021 or
that Entertainment dipped in mid-2019.
(3) Granularity therefore sets the question. Use monthly rows to attribute peaks to events and yearly rows to compare
eras. The band height scales roughly with the bin length (55, 150, 469 posts), which is why a common colour does not mean
a common count \emph{across} panels.

\subsection{Variation 3: 50 subreddits, score anomaly}
\figfull{Fig_39_horizon_subreddits_monthly.png}{Horizon plot of the monthly median $\log_{10}$(score) of each subreddit's top
posts relative to the subreddit's overall median, for all 50 subreddits (rows grouped by topic, label colour = topic).
$h=0.15$ dex (a factor of 1.41); deviations beyond 0.45 dex saturate. Months without a top post are blank.}{fig:39}

\textbf{Inferences.}
(1) Fifty series fit on one page only because each row is about 4\,mm high. A horizon plot is the only practical way to
show this many series at monthly resolution.
(2) The share of subreddits whose yearly median is above their own median peaks in 2017 (60\,\%) and 2020 (64\,\%) and falls
to 18\,\% in 2023 and 16\,\% in 2024. The most recent top posts are both fewer \emph{and} lower-scoring, again an effect of
survivorship plus declining engagement.
(3) Individual communities stand out. Art's 2023--24 top posts are a full order of magnitude below its usual score (median
deviation $-1.08$ dex), and GetMotivated, askscience, personalfinance and Showerthoughts are 0.3--0.5 dex below theirs.
These communities lost reach after 2022. AskReddit and travel were above their medians in 2016.
(4) Blank stretches carry information too: 18\,\% of subreddit-months have no top post. memes (53\,\%) and
relationship\_advice (52\,\%) have the most gaps because their top posts are concentrated in a few years. A missing month is
a gap, not ``average'', and the plot deliberately leaves it white.

\subsection{Ridgeline plots}
Each ridge is a Gaussian kernel density estimate~\cite{silverman1986}
\begin{equation}
\hat f_g(x)=\frac{1}{n_g h}\sum_{i\in g}\phi\!\left(\frac{x-x_i}{h}\right),
\end{equation}
with \textbf{one fixed bandwidth for all groups} ($h=0.06$ dex for $\log_{10}$ score and $h=0.08$ for $\log_{10}$(comments+1)).
Shapes are therefore comparable, which would not be the case with a separate Scott or Silverman bandwidth per group. All
groups share the $x$-axis. Heights are scaled by the global maximum density with overlap factor 1.6--2.2, and a white tick
marks each median. We show three variations with different groupings and orderings.

\figfullw{Fig_40_ridgeline_subreddits_score.png}{Ridgeline of $\log_{10}$(score) for the ~1,000 top posts of each of the 50
subreddits, ordered by median (highest at top), colour = topic, bandwidth 0.06.}{fig:40}{0.75}

\textbf{Inferences (Fig.~\ref{fig:40}).}
(1) Medians span 1.9 orders of magnitude, from memes ($10^{5.13}\approx135$k) and pics (125k) down to listentothis
($10^{3.24}\approx1.7$k). Being a ``top post'' means very different things in different communities.
(2) The top of the ranking is dominated by image and meme communities (memes, pics, funny, aww, interestingasfuck, gaming,
gifs), the middle by news and discussion, and the bottom by niche-interest communities (history, travel,
InternetIsBeautiful, personalfinance, askscience, Documentaries), which matches the topic colours.
(3) Shapes differ, not just positions. The largest default communities have extremely narrow ridges: todayilearned, aww and
memes have a standard deviation of only 0.07 dex, so all their top 1,000 posts lie within a factor of about 1.4 of each
other. AskReddit is similarly tight, with 80\,\% of its top posts between $10^{4.76}$ and $10^{4.96}$. At the other end, Art
(0.49 dex), relationship\_advice (0.26) and InternetIsBeautiful (0.25) have wide, skewed ridges, where a few viral posts
sit far above a modest bulk.

\figfullw{Fig_41_ridgeline_years_score.png}{Ridgeline of $\log_{10}$(score) of all top posts by posting year, 2015--2024,
ordered by time, bandwidth 0.06 (2011--14 omitted: 176 posts).}{fig:41}{0.75}

\textbf{Inferences (Fig.~\ref{fig:41}).}
(1) The median score rises from $10^{4.07}$ (11.7k, 2015) to a maximum of $10^{4.68}$ (47.4k, 2020) and then falls to
$10^{4.23}$ (17.1k, 2024). The distributions shift right until 2020 and left after it, the same arc as the horizon
counts.
(2) The 2020, 2021 and 2022 ridges are almost identical (standard deviation 0.41--0.42 dex, 10th percentile
$10^{4.02}$--$10^{4.03}$). This is the plateau of engagement in the corpus. 2015--16 have the same width but are shifted
left by about 0.5 dex, because their communities were smaller when they were posted.
(3) 2023 and especially 2024 are wider (standard deviation 0.52 and 0.56) and the 2024 ridge is \emph{bimodal}: a low mode at
$10^{3.3}$--$10^{3.4}$ (about 2k) next to the usual mode at $10^{4.1}$--$10^{4.6}$. Recent posts in niche communities reach
the top 1,000 with few votes, while mass-audience recent posts have not yet accumulated theirs. This is a censoring effect
the reader must keep in mind.

\figfullw{Fig_42_ridgeline_topics_comments.png}{Ridgeline of $\log_{10}$(comments+1) per post by topic category, ordered by
median, bandwidth 0.08 (all 49,266 posts).}{fig:42}{0.75}

\textbf{Inferences (Fig.~\ref{fig:42}).}
(1) News, science and tech generates the most comments per top post (median 2,145), then Q\&A (1,263), Entertainment (1,201),
Memes (1,136) and Photos (1,105). Creative and lifestyle is far behind (373), because people upvote art, food and DIY but
rarely debate it.
(2) Q\&A has the widest ridge (standard deviation 0.53 dex): AskReddit threads collect 6k--34k comments (10th to 90th
percentile), while advice posts in LifeProTips or personalfinance receive hundreds. Its right tail is the longest of all
topics, with a 99th percentile of $10^{4.57}\approx37$k comments.
(3) Compared with Fig.~\ref{fig:40}, where Memes and Photos lead on score, the comment ranking reverses the visual
communities. Score and discussion are different kinds of engagement, consistent with the negative ratio--comments relation
found by brushing in I2a.

\subsection{Observations and limitations (I3)}
Horizon plots compress many series about 4$\times$ more efficiently than line charts and make above/below-baseline phases
pre-attentive. Granularity and baseline must be chosen to match the question. Ridgelines with a fixed bandwidth allow
honest comparison of distribution shape and location across 50 groups. Limitations: the counts reflect survivorship in an
all-time-top sample, not posting volume. KDEs smooth the hard lower bound of the corpus (each subreddit's top-1,000
threshold), so the left tails are slightly widened. Horizon plots need a legend to be read correctly.
